% Generated draft: requires mathematical and source review.
\documentclass[11pt]{article}
\usepackage[margin=25mm]{geometry}
\usepackage{fontspec}
\setmainfont{Latin Modern Roman}
\usepackage{amsmath,amssymb,mathtools,unicode-math}
\setmathfont{Latin Modern Math}
\usepackage{enumitem,needspace,xurl,hyperref}
\setlength{\parindent}{0pt}
\setlength{\parskip}{5pt}
\setlength{\emergencystretch}{3em}
\newcommand{\N}{\mathbb{N}}
\newcommand{\Z}{\mathbb{Z}}
\newcommand{\Q}{\mathbb{Q}}
\newcommand{\R}{\mathbb{R}}
\newcommand{\C}{\mathbb{C}}
\newcommand{\F}{\mathbb{F}}
\newcommand{\A}{\mathbb{A}}
\newcommand{\PP}{\mathbb{P}}
\newcommand{\E}{\mathbb{E}}
\newcommand{\eps}{\varepsilon}
\newcommand{\dd}{\,\mathrm{d}}
\begin{document}
\section*{UnsolvedMath ID: 701683}
\textbf{Source problem code:} catalogue entry
\subsection{Definitions and mathematical statement}
For each positive integer \(k\), let \(n=k^2\), and let \(S_n\) be the symmetric group on \(n\) letters. Choose \(\sigma_k\in S_{k^2}\) whose disjoint cycle lengths are exactly \(1,3,5,\ldots,2k-1\), each occurring once; these lengths sum to \(k^2\). Interpret the character table as the ordinary complex irreducible character table. Its rows are indexed by partitions \(\lambda\vdash k^2\); write \(\chi^\lambda\) for the corresponding character and \(p(m)\) for the number of partitions of \(m\). The values \(\chi^\lambda(\sigma_k)\) are integers. An entry is odd if it belongs to \(2\Z+1\), including negative odd integers. Define
\[
 a_k=\frac{\#\{\lambda\vdash k^2:\chi^\lambda(\sigma_k)\in 2\Z+1\}}{p(k^2)}.
\]
This proportion is independent of the choice of representative \(\sigma_k\) of the specified conjugacy class. Determine
\[
\lim_{k\to\infty}a_k,
\]
including establishing its existence. The source asks for a limiting proportion without separately asserting that the limit exists; the ordinary complex character-table interpretation is made explicit here.
\subsection{Short English statement}
For permutations in \(S_{k^2}\) with one cycle of each odd length from \(1\) to \(2k-1\), determine the limiting fraction of irreducible character values that are odd integers as \(k\to\infty\). This is the problem attributed to D. Gluck; the supplied registry labels it open.
\subsection{Sources}
\begin{itemize}
\item Supplied local UnsolvedMath registry record, permanent ID \(701683\), exact source problem code \texttt{catalogue entry}, titled \emph{Kourovka Notebook Problem 20.25}. Its supplied status is \texttt{open}; this label is not an independently verified literature assessment.
\item \emph{The Kourovka Notebook (Unsolved Problems in Group Theory)}, original problem number 20.25, proposed by D. Gluck in 2022. The supplied provenance identifies the 20th issue (2022) and the 21st issue edition (2026), arXiv:1401.0300v48. The supplied collection description names E. I. Khukhro and V. D. Mazurov as editors.
\item Background reference supplied in the statement: D. Gluck, \emph{Proc. Amer. Math. Soc.} \textbf{147} (2019), 1005--1011, specifically pages 1005--1006. The article title and a direct link are not supplied.
\end{itemize}
\end{document}
